\documentclass[12pt,a4paper]{article}
\usepackage{ugmath}
\usepackage{float}
\usepackage{placeins}
\usepackage{array}
\usepackage[labelfont=bf,labelsep=space]{caption}
\newcommand{\paren}[1]{\left( #1 \right)}

\title{Solucion de Problemas\\
Taller de calculo}

\author{Equipo 5}

\begin{document}
    \maketitle

    \textbf{Lema 1. (Criterio de comparación)} Sean $(a_{n})$ y $(b_{n})$ sucesiones tales que
    $0\leq a_{n}\leq b_{n}$ para todo $n\in\N$ si $\sum b_{n}$ converge
    entonces $\sum a_{n}$ converge

    \begin{proof}
        Sean
        \begin{equation*}
            A_{n}=\sum_{k=1}^{n}a_{k}\quad\text{ y }\quad B_{n}=\sum_{k=1}^{n}b_{k}
        \end{equation*}
        las sucesiones de sumas parciales. Como $a_{n}\geq0$ para todo $n$,
        \begin{equation*}
            A_{n+1}=A_{n}+a_{n+1}\geq A_{n}.
        \end{equation*}
        Como $a_{k}\leq b_{k}$ para todo $k$,
        \begin{equation*}
            A_{n}=\sum_{k=1}^{n}a_{k}\leq \sum_{k=1}^{n}b_{k}=B_{n}.
        \end{equation*}
        Por hipotesis $\sum_{k=1}^{\infty}b_{k}$ converge, asi que $(B_{n})$
        converge a algún $L$. Como $b_{k}\geq0$, $(B_{n})$ tambien es creciente.
        Por lo tanto $B_{n}\leq L$ para todo $n$. Combinando ambas desigualdades
        \begin{equation*}
            A_{n}\leq B_{n}\leq L.\quad\forall n.
        \end{equation*}
        Así, $(A_{n})$ es creciente y acotada superiormente por $L$. Por el teorema de
        Weierstrass $(A_{n})$ converge, es decir,
        \begin{equation*}
            \lim_{n\to\infty}A_{n}\text{ existe.}
        \end{equation*}
        Esto significa que la serie
        \begin{equation*}
            \sum_{k=1}^{\infty}a_{k}
        \end{equation*}
        converge.
    \end{proof}

    \begin{ejercicio}
        Prueba el Teorema 1.
    \end{ejercicio}

    \begin{proof}
        Sea
        \begin{equation*}
            S_{n}=\sum_{i=1}^{n}x_{i}
        \end{equation*}
        la $n-$ésima suma parcial de la serie, y sea $S$ el valor al que converge, es
        \begin{equation*}
            \lim_{n\to\infty}S_{n}=S.
        \end{equation*}
        Por la definición 2, esto significa que
        \begin{equation}
            \forall\e\gt0\exists n_{0}\in\N\text{ tal que }n\geq n_{0}\implies\abs{S_{n}-S}\lt\frac{\e}{2}.
        \end{equation}
        Nótese que para $n\geq2$,
        \begin{equation}
            x_{n}=S_{n}-S_{n-1}.
        \end{equation}
        Sea $\e\gt0$ dado. Por (1), existe $n_{0}\in\N$ tal que
        \begin{equation*}
            n\geq n_{0}\implies\abs{S_{n}-S}\lt\frac{\e}{2}.
        \end{equation*}
        Sea $n\geq n_{0}+1$ arbitrario. Entonces tanto $n$ como $n-1$ son
        mayores o iguales que $n_{0}$, así que se cumplen simultaneamente
        \begin{equation*}
            \abs{S_{n}-S}\lt\frac{\e}{2}\quad\text{ y }\quad\abs{S_{n-1}-S}\lt\frac{\e}{2}.
        \end{equation*}
        Usando (2) y la desigualdad del triangulo
        \begin{equation*}
            \abs{x_{n}}=\abs{(S_{n}-S)-(S_{n-1}-S)}\leq\abs{S_{n}-S}+\abs{S_{n-1}-S}\lt\frac{\e}{2}+\frac{\e}{2}=\e.
        \end{equation*}
        Hemos demostrado que, dado cualquier $\e\gt0$, tomando $n_{1}:=n_{0}+1$ se cumple,
        \begin{equation*}
            n\geq n_{1}\implies \abs{x_{n}-0}\lt\e.
        \end{equation*}
        Por esto precisamente, por la definición 2, la afirmación de que
        \begin{equation*}
            \lim_{n\to\infty}x_{n}=0.
        \end{equation*}
    \end{proof}

    \begin{ejercicio}
        Usando la Definición 2 prueba que la serie
        \begin{equation*}
            \sum_{n=1}^{\infty}\frac{1}{2^{n}}
        \end{equation*}
        es convergente.
    \end{ejercicio}

    \begin{proof}
        Sea
        \begin{equation*}
            S_{n}=\sum_{k=1}^{n}\frac{1}{2^{k}}
        \end{equation*}
        la sucesión de sumas parciales de la serie. Por la fórmula
        de la suma de una progesión geométrica,
        \begin{equation*}
            S_{n}=\frac{\frac12\paren{1-\paren{\frac12}^{n}}}{1-\frac{1}{2}}=1-\frac{1}{2^{n}}.
        \end{equation*}
        Afirmamos que
        \begin{equation*}
            \lim_{n\to\infty}S_{n}=1.
        \end{equation*}
        En efecto, sea $\e\gt0$. Como $\e\gt0$, por la propiedad arquimediana
        existe $N\in\N$ tal que
        \begin{equation*}
            2^{N}\gt\frac{1}{\e}.
        \end{equation*}
        Equivalentemente
        \begin{equation*}
            \frac{1}{2^{N}}\gt\e.
        \end{equation*}
        Ahora, si $n\geq N$, se tiene
        \begin{equation*}
            2^{n}\geq 2^{N},
        \end{equation*}
        y, como la función $x\mapsto\frac{1}{x}$ es estrictamente decreciente en
        $(0,\infty)$,
        \begin{equation*}
            \frac{1}{2^{n}}\leq\frac{1}{2^{N}}\lt\e.
        \end{equation*}
        Además,
        \begin{equation*}
            \abs{S_{n}-1}=\abs{1-\frac{1}{2^{n}}-1}=\frac{1}{2^{n}}.
        \end{equation*}
        Por consiguiente,
        \begin{equation*}
            \abs{S_{n}-1}-\frac{1}{2^{n}}\leq\frac{1}{2^{N}}\lt\e.
        \end{equation*}
        Como $\e\gt0$ fue arbitrario
        \begin{equation*}
            \lim_{n\to\infty}S_{n}=1.
        \end{equation*}
        Por definición de convergencia de sieries,
        \begin{equation*}
            \sum_{n=1}^{\infty}\frac{1}{2^{n}}
        \end{equation*}
        converge y su suma es 1.
    \end{proof}

    \begin{ejercicio}[\textbf{(Criterio de Raabe)}]
        Supón que $a_{n}\gt0$ para todo $n$ y que
        \begin{equation*}
            \lim_{n\to\infty}n\paren{\frac{a_{n}}{a_{n+1}}-1}=\ell,
        \end{equation*}
        entonces,
        \begin{itemize}
            \item si $\ell\gt1$ la serie $\sum a_{n}$ converge;
            \item si $\ell\lt1$ la serie diverge.
        \end{itemize}
    \end{ejercicio}

    \begin{proof}
        Caso $\ell\gt1$.\\
        Como $\ell\gt1$, podemos elegir $\e=\frac{\ell-1}{2}\gt0$. Por la
        definición de limite, existe $N\in\N$ tal que para todo $n\geq N$:
        \begin{equation*}
            \abs{n\paren{\frac{a_{n}}{a_{n-1}}-1}-\ell}\lt\e.
        \end{equation*}
        Abriendo el valor absoluto
        \begin{equation*}
            -\e\lt n\paren{\frac{a_{n}}{a_{n+1}}-1}-\ell\lt\e.
        \end{equation*}
        Sumando $\ell$ en ambos lados,
        \begin{equation*}
            \e-\ell\lt n\paren{\frac{a_{n}}{a_{n+1}}-1}\lt\ell+\e
        \end{equation*}
        llamamos $r=\ell-\e=\frac{\ell+1}{2}$, que satisface $r\gt1$ porque
        $\ell\gt1$. De la desigualdad izquierda obtenemos que para todo
        $n\geq N$
        \begin{equation*}
            n\paren{\frac{a_{n}}{a_{n+1}}-1}\geq r.
        \end{equation*}
        Multiplicando por $a_{n+1}\gt0$,
        \begin{equation*}
            na_{n}-na_{n+1}\geq ra_{n+1}
        \end{equation*}
        sumando y restando $a_{n+1}$ del lado izquierdo obtenemos
        \begin{align*}
            na_{n}-na_{n+1}+a_{n+1}-a_{n+1}&\geq ra_{n+1}\\
            na_{n}-na_{n+1}-a_{n+1}&\geq ra_{n+1}-a_{n+1}\\
            na_{n}-(n+1)a_{n+1}&\geq(r-1)a_{n+1}
        \end{align*}
        Definamos $b_{n}=na_{n}$. La desigualdad anterior dice
        \begin{equation*}
            b_{n}-b_{n+1}\geq(r-1)a_{n+1},
        \end{equation*}
        es decir, $b_{n}$ es estrictamente decreciente para $n\geq N$.
        Sumando telescopicamente de $N$ a $M$
        \begin{equation*}
            \sum_{n=N}^{M}(b_{n}-b_{n+1})=b_{N}-b_{M+1}\geq(r-1)\sum_{n=N}^{M}a_{n+1}.
        \end{equation*}
        Como $b_{M+1}=(M+1)a_{M+1}\gt0$, se tiene $b_{N}-b_{M+1}\leq b_{N}$,
        y, entonces
        \begin{equation*}
            (r-1)\sum_{n=N}^{M}a_{n+1}\leq b_{N}.
        \end{equation*}
        Dividiendo entre $(r-1)\gt0$,
        \begin{equation*}
            \sum_{n=N}^{M}a_{n+1}\leq\frac{b_{N}}{r-1}.
        \end{equation*}
        Consideramos las sumas parciales
        \begin{equation*}
            S_{M}=\sum_{n=1}^{M}a_{n}
        \end{equation*}
        y las partimos en dos trozos
        \begin{equation*}
            S_{M}=\sum_{n=1}^{N}a_{n}+\sum_{n=N+1}^{M}a_{n}
        \end{equation*}
        y
        \begin{equation*}
            \sum_{n=N+1}^{M}a_{n}=\sum_{n=N}^{M-1}a_{n+1}\leq\frac{b_{N}}{r-1}.
        \end{equation*}
        Entonces para todo $M$,
        \begin{equation*}
            S_{M}\leq\sum_{n=1}^{N}a_{n}+\frac{b_{N}}{r-1}:=K
        \end{equation*}
        donde $K$ no depende de $M$. Como $S_{M}$ es creciente y acotada
        superiormente por $K$, por el teorema de Weierstrass, $(S_{M})$. Por lo tanto
        $\sum a_{n}$ converge.
    \end{proof}

    \begin{ejercicio}[\textbf{(Criterio de condensación de Cauchy)}]
        Si la sucesión $(a_{n})_{n\geq1}$ satisface $a_{n}\geq a_{n+1}\geq0$
        para todo $n\in\N$, entonces $\sum a_{n}$ converge si y solo si $\sum2^{k}a_{2^{k}}$
        converge.
    \end{ejercicio}

    \begin{proof}
        Para cada $k\in\N$, consideremos el bloque de indices $\{2^{k},2^{k}+1,\dots,2^{k+1}-1\}$,
        el cual contiene exactamente $2^{k}$ elementos. Como $(a_{n})$ es no creciente, para todo $i$
        en dicho bloque se cumple $a_{2^{k+1}}\leq a_{i}\leq a_{2^{k}}$. Sumando sobre los $2^{k}$ indices
        del bloque
        \begin{equation*}
            2^{k}a_{2^{k+1}}\leq\sum_{i=2^{k}}^{2^{k+1}-1}a_{i}\leq2^{k}a_{2^{k}}.
        \end{equation*}
        Sumando ahora sobre los bloques $k=1,\dots,m$
        \setcounter{equation}{0}
        \begin{equation}
            \sum_{k=1}^{m}2^{k}a_{2^{k+1}}\leq\sum_{k=1}^{m}\sum_{i=2^{k}}^{2^{k+1}-1}a_{i}\leq\sum_{k=1}^{m}2^{k}a_{2^{k}}.
        \end{equation}
        Simplifiquemos cada extremo de (1). Por un lado,
        \begin{equation*}
            \sum_{k=1}^{m}2^{k}a_{2^{k+1}}=\frac12\sum_{k=1}^{m}2^{k+1}a_{2^{k+1}}=\frac12\paren{\sum_{k=1}^{m+1}2^{k}a_{2^{k}}-2a_{2}}.
        \end{equation*}
        Por otro lado, los bloques disjuntos y su union es $\{2,3,\dots,2^{m+1}-1\}$, asi que
        \begin{equation*}
            \sum_{k=1}^{m}\sum_{i=2^{k}}^{2^{k+1}-1}a_{i}=\sum_{i=1}^{2^{m+1}-1}a_{i}-a_{1}.
        \end{equation*}
        Sustituyendo ambas simplificaciones en (1) obtenemos
        \begin{equation}
            \frac12\paren{\sum_{k=1}^{m+1}2^{k}a_{2^{k}}-2a_{2}}\leq\sum_{i=1}^{2^{m+1}-1}a_{i}-a_{1}\leq\sum_{k=1}^{m}2^{k}a_{2^{k}}
        \end{equation}
        Con esta desigualdad podemos probar ambas implicaciones.\\
        Supongamos primero que $\sum2^{k}a_{2^{k}}$ converge. Entonces sus sumas parciales estan acotadas,
        y por la desigualdad derecha de (2), las sumas parciales $\sum_{i=1}^{2^{m+1}-1}a_{i}$ tambien estan acotadas superiormente.
        Como $a_{i}\geq0$, las sumas parciales de $\sum a_{n}$ forman una sucesión no decreciente y acotada, luego convergen. Así,
        $\sum a_{n}$ converge.\\
        Supongamos ahora que $\sum a_{n}$ converge. Entonces, en particular,
        \begin{equation*}
            \lim_{m\to\infty}\sum_{i=1}^{2^{m+1}-1}a_{i}
        \end{equation*}
        existe, y por la desigualdad izquierda de (2),
        \begin{equation*}
            \sum_{k=1}^{m+1}2^{k}a_{2^{k}}\leq2\sum_{i=1}^{2^{m+1}-1}a_{i}-2a_{1}+2a_{2},
        \end{equation*}
        de modo que las sumas parciales $\sum_{k=1}^{m+1}2^{k}a_{2^{k}}$ esta acotadas superiormente. Como
        $2^{k}a_{2^{k}}\geq0$, forman una sucesión no decreciente y acotada, luego convergen. Así, $\sum2^{n}a_{2^{n}}$
        converge.\\
        De ambas implicaciones concluimos que
        \begin{equation*}
            \sum_{n=1}^{\infty}a_{n}\text{ converge}\ssi\sum_{n=1}^{\infty}2^{n}a_{2^{n}}\text{ converge}
        \end{equation*}
    \end{proof}

    \begin{ejercicio}
        Si la serie $\sum\abs{x_{n}}$ converge entonces $\sum x_{n}$
        converge.
    \end{ejercicio}

    \begin{proof}
        Para cada $n$, define
        \begin{equation*}
            y_{n}=x_{n}+\abs{x_{n}}.
        \end{equation*}
        Como $x_{n}\leq\abs{x_{n}}$ siempre, se tiene $y\geq0$.
        Y como $x_{n}\geq-\abs{x_{n}}$, tambien se tiene $y_{n}\leq2\abs{x_[n]}$.
        Entonces
        \begin{equation*}
            0\leq y_{n}\leq 2\abs{x_{n}}
        \end{equation*}
        para todo $n$. Por hipotesis $\sum\abs{x_[n]}$ converge, así que $\sum2\abs{x_{n}}$
        también converge. Por el crierio de comparación la serie
        \begin{equation*}
            \sum y_{n}
        \end{equation*}
        converge. Ahora despejamos $x_{n}$,
        \begin{equation*}
            x_{n}=y_{n}-\abs{x_{n}}.
        \end{equation*}
        La diferencia de series convergentes converge. Por lo tanto
        $\sum x_{n}$ converge.
    \end{proof}

    \begin{ejercicio}
        Prueba que la serie
        \begin{equation*}
            \sum_{n=1}^{\infty}(-1)^{n+1}\frac1n
        \end{equation*}
        converge.
    \end{ejercicio}

    \begin{proof}
        Primero veamos que la subsucesion de sumas parciales pares, $S_{2n}$,
        es creciente. Calculamos
        \begin{equation*}
            S_{2n+2}-S_{2n}=\frac{1}{2n+1}-\frac{1}{2n+2}.
        \end{equation*}
        Como $2n+1\lt2n+2$, se tiene
        \begin{equation*}
            \frac{1}{2n+1}\gt\frac{1}{2n+2},
        \end{equation*}
        asi que
        \begin{equation*}
            s_{2n+2}-S_{2n}\gt0\implies S_{2n+2}\gt S_{2n}.
        \end{equation*}
        Entonces $(S_{2n})$ es estrictamente creciente. De forma analoga
        \begin{equation*}
            S_{2n+3}-S_{2n+1}=-\frac{1}{2n+2}+\frac{1}{2n+3}=\frac{1}{2n+3}-\frac{1}{2n+2}\lt0
        \end{equation*}
        porque $2n+3\gt 2n+2$. Entonces $(S_{2n+1})$ es estrictamente decreciente.\\
        Toda suma par es menor que toda suma impar. En efecto, primero para un mismo bloque
        \begin{equation*}
            S_{2n+1}-S_{2n}=\frac{1}{2n+1}\gt0\implies S_{2n}\lt S_{2n+1},
        \end{equation*} 
        como $(S_{n})$ crece y $(S_{2n+1})$ decrece, y $S_{2n}\lt S_{2n+1}$
        en cada paso, se concluye que para tods $n,m$
        \begin{equation*}
            S_{2n}\lt S_{2m+1}.
        \end{equation*}
        Ahora, $(S_{2n})$ es creciente y acotada superiormente, por el teorema de Weierstrass
        converge, es decir
        \begin{equation*}
            S_{2n}\to L_{1}.
        \end{equation*}
        $(S_{2n+1})$ es decreciente y acotada inferiormente, por lo tanto, tambien converge,
        es decir,
        \begin{equation*}
            S_{2n+1}\to L_{2}.
        \end{equation*}
        Como $S_{2n+1}-S_{2n}=\frac{1}{2n+1}\to0$ cuando $n\to\infty$,
        tomando el limite en ambos lados
        \begin{equation*}
            L_{2}-L_{1}=\lim_{n\to\infty}\frac{1}{2n+1}=0\implies L_{1}=L_{2}=L.
        \end{equation*}
        Ya sabemos que las sumas pares y las impares convergen al mismo limite $L$.
        Como toda la sucesión $(S_{n})$ esta formada exactamente por estas dos subsucesiones,
        se concluye que
        \begin{equation*}
            S_{n}\to L.
        \end{equation*}
        Por lo tanto la serie
        \begin{equation*}
            \sum_{k=1}^{\infty}(-1)^{n+1}\frac{1}{n}
        \end{equation*}
        converge.
    \end{proof}

    \begin{ejercicio}
        Sea $(x_{n})_{n\geq1}$ una sucesión decreciente de numeros positivos tal que
        $\lim_{n\to\infty}x_{n}=0$, entonces
        \begin{equation*}
            \sum_{n=1}^{\infty}(-1)^{n+1}x_{n}
        \end{equation*}
        es converge. Mas aun, si
        \begin{equation*}
            S_{k}=\sum_{n=1}^{k}(-1)^{n+1}x_{n}\quad\text{ y }\quad S=\sum_{n=1}^{\infty}(-1)^{n+1}x_{n}
        \end{equation*}
        entonces
        \begin{equation*}
            \abs{S-S_{k}}\leq x_{k+1}\quad\text{para todo }k\geq1.
        \end{equation*}
    \end{ejercicio}

    \begin{proof}
        Demostremos primero la convergencia. Consideremos las subsucesiones de sumas parciales con
        índice par e índice impar, $S_{2k}$ y $S_{2k+1}$. Notemos primero que
        \begin{equation*}
            S_k = x_1-x_2+x_3-x_4+\cdots+(-1)^{k+1}x_k.
        \end{equation*}

        Si $k$ es impar, podemos agrupar
        \begin{equation*}
            S_k = x_1-(x_2-x_3)-(x_4-x_5)-\cdots-(x_{k-1}-x_k),
        \end{equation*}
        y como $(x_n)$ es no creciente, cada paréntesis es no negativo, de donde $S_k\le x_1$. Si $k$
        es par, análogamente
        \begin{equation*}
            S_k = x_1-(x_2-x_3)-(x_4-x_5)-\cdots-(x_{k-2}-x_{k-1})-x_k \le x_1.
        \end{equation*}
        En cualquier caso, $S_k\le x_1$; es decir, $(S_k)_{k\ge1}$ está acotada superiormente por $x_1$.

        Ahora bien, para la subsucesión de índices pares,
        \begin{equation*}
            S_{2k+2}-S_{2k}=x_{2k+1}-x_{2k+2}\ge 0,
        \end{equation*}
        pues $(x_n)$ es no creciente. Así, $(S_{2k})_{k\ge1}$ es creciente, y por lo anterior está
        acotada superiormente; por el teorema de sucesiones monótonas acotadas, $(S_{2k})_{k\ge1}$
        converge, digamos a $S$.

        Por otro lado, para todo $k\in\mathbb{N}$ se cumple la identidad
        \begin{equation*}
            S_{2k+1}=S_{2k}+x_{2k+1}.
        \end{equation*}
        Como $(S_{2k})$ converge a $S$ y, por hipótesis, $\lim_{n\to\infty}x_n=0$ (en particular
        $\lim_{k\to\infty}x_{2k+1}=0$), se sigue que
        \begin{equation*}
            \lim_{k\to\infty} S_{2k+1}=\lim_{k\to\infty}S_{2k}+\lim_{k\to\infty}x_{2k+1}=S+0=S.
        \end{equation*}

        Como ambas subsucesiones —la de índices pares y la de índices impares, que en conjunto
        cubren todo $\mathbb{N}$— convergen al mismo límite $S$, concluimos que $(S_k)_{k\ge1}$
        converge a $S$, es decir, la serie $\sum(-1)^{n+1}x_n$ converge.

        Demostremos ahora la cota del resto $|S-S_k|\le x_{k+1}$. Fijemos $k\ge1$ y escribamos la
        cola de la serie a partir del índice $k+1$:
        \begin{equation*}
            R_k:=S-S_k=\sum_{i=k+1}^{\infty}(-1)^{i+1}x_i.
        \end{equation*}

        Factoricemos el signo global escribiendo $R_k=(-1)^{k+2}T_k$, donde
        \begin{equation*}
            T_k := x_{k+1}-x_{k+2}+x_{k+3}-x_{k+4}+\cdots
        \end{equation*}

        Agrupando por parejas a partir del segundo término,
        \begin{equation*}
            T_k = x_{k+1}-(x_{k+2}-x_{k+3})-(x_{k+4}-x_{k+5})-\cdots,
        \end{equation*}
        y como $(x_n)$ es no creciente, cada paréntesis es no negativo, de donde $T_k\le x_{k+1}$.

        Agrupando en cambio desde el primer término,
        \begin{equation*}
            T_k = (x_{k+1}-x_{k+2})+(x_{k+3}-x_{k+4})+\cdots \ge 0,
        \end{equation*}
        pues cada paréntesis es no negativo.

        Por lo tanto $0\le T_k\le x_{k+1}$, y como $R_k=\pm T_k$ según la paridad de $k$, se tiene
        \begin{equation*}
            |S-S_k|=|R_k|=T_k\le x_{k+1}.
        \end{equation*}

        Como $k\ge1$ era arbitrario, la desigualdad se cumple para todo $k\ge1$.
    \end{proof}

    \begin{ejercicio}
        Prueba que si la serie $\sum x_{n}$, es absolutamente
        convergente entonces todos sus reordenamientos convergen al mismo valor.
    \end{ejercicio}

    \begin{proof}
        Llamemos
        \begin{equation*}
            S_{n}=\sum_{k=1}^{n}x_{k}\quad\text{ y }\quad T_{n}=\sum_{k=1}^{n}x_{\sigma(k)}
        \end{equation*}
        a las sumas parciales de la serie original y del reordenamiento respectivamente.
        Sea $\e\gt0$. Como $\sum \abs{x_{n}}$ converge, sus sumas parciales forman una sucesión
        de cauchy, lo que en particular significa que existe $N$ tal que
        \begin{equation*}
            \sum_{n=N+1}^{\infty}\abs{x_{n}}\lt\e
        \end{equation*}
        Ahora, como $\sigma$ es una biyección, los índices $1,2,\dots,N$ aparecen en todos los
        lugares de la lista $\sigma(1),\sigma(2),\sigma(3),\dots$. Entonces existe $M$ tal que
        \begin{equation*}
            \{1,2,\dots,N\}\subseteq\{\sigma(1),\sigma(2),\dots,\sigma(M)\}.
        \end{equation*}
        Es decir, para $m\geq M$, los terminos $x_{1},x_{2},\dots,x_{N}$
        ya estan dentro de $T_{m}$. Acotamos ahora $\abs{T_{m}-S_{N}}$ para
        $m\geq M$. Al restar $S_N$ a $T_m$, los términos $x_1,\dots,x_N$ se cancelan
        (pues, como $m\geq M$, todos ellos aparecen entre $x_{\sigma(1)},\dots,x_{\sigma(m)}$).
        Lo que queda es una suma de algunos de los términos $x_{\sigma(k)}$ con $\sigma(k)>N$;
        como $\sigma$ es inyectiva, estos son términos $x_{n}$ con $n>N$ y todos distintos
        entre sí. Es decir,
        \begin{equation*}
            T_{m}-S_{N}=x_{n_{1}}+x_{n_{2}}+\cdots+x_{n_{r}}
        \end{equation*}
        para ciertos índices distintos $n_{1},\dots,n_{r}>N$. Por la desigualdad del triángulo,
        \begin{equation*}
            \abs{T_{m}-S_{N}}\leq\abs{x_{n_{1}}}+\abs{x_{n_{2}}}+\cdots+\abs{x_{n_{r}}}.
        \end{equation*}
        Como los índices $n_{1},\dots,n_{r}$ son distintos y todos mayores que $N$, esta suma es
        parte de la suma de la cola $\sum_{n=N+1}^{\infty}\abs{x_{n}}$; y como esta última tiene
        solo términos no negativos, cualquier suma parcial de sus términos está acotada por la suma
        total. Por lo tanto,
        \begin{equation*}
            \abs{T_{m}-S_{N}}\leq\sum_{n=N+1}^{\infty}\abs{x_{n}}\lt\e.
        \end{equation*}
        Acotamos $\abs{S-S_{N}}$
        \begin{equation*}
            \abs{S-S_{N}}=\abs{\sum_{n=N+1}^{\infty}x_{n}}\leq \sum_{n=N+1}^{\infty}\abs{x_{n}}\lt\e
        \end{equation*}
        Juntamos todo con la desigualdad del triangulo. Para $m\geq M$
        \begin{equation*}
            \abs{T_{m}-S}\leq\abs{T_{m}-S_{N}}+\abs{S_{N}-S}\lt \e+\e=2\e
        \end{equation*}
        Como $\e\gt0$ era arbitrario, esto prueba que
        \begin{equation*}
            T_{m}\to S.
        \end{equation*}
    \end{proof}

    \begin{ejercicio}
        Supon ahora que la serie $\sum x_{n}$ es condicionalmente
        convergente. Sean $p_{1},p_{2},p_{3},\dots$ los términos
        estrictamente positivos en la sucesión $(x_{n})_{n\geq1}$
        en el mismo orden en el que aparecen en la serie original;
        de la misma manera, sea $-q_{1},-q_{2},-q_{3},\dots$ la sucesión
        de terminos estrictamente negativos, en el mismo orden en que
        aparecen en $\{x_{n}\}$. Prueba que ambas sucesiones $(p_{n})_{n\geq1}$
        y $(q_{n})_{n\geq1}$ tienen una infinidad de término y que ambas series $\sum p_{n}$
        y $\sum q_{n}$ divergen a $+\infty$
    \end{ejercicio}

    \begin{proof}
        Notemos primero que, como $p_n\ge0$ y los términos nulos no contribuyen a la suma, la serie
        $\sum_n p_n$ converge si y solo si converge la serie formada exclusivamente por sus términos
        estrictamente positivos, $\sum_k p_k'$, y en tal caso ambas tienen la misma suma
        (análogamente para $q_n$ y $q_k'$). Basta entonces con estudiar la convergencia de
        $\sum p_n$ y $\sum q_n$.

        Supongamos, buscando una contradicción, que $\sum p_n$ converge. Como por hipótesis
        $\sum x_n$ converge, la serie $\sum(x_n-p_n)$ converge, al ser diferencia de dos series
        convergentes. Pero $x_n-p_n=-q_n$, así que $\sum(-q_n)$ converge, y por lo tanto
        $\sum q_n=-\sum(-q_n)$ también converge.

        Pero entonces, como $|x_n|=p_n+q_n$, la serie $\sum|x_n|=\sum p_n+\sum q_n$ sería suma de
        dos series convergentes y, por lo tanto, convergería. Esto contradice que $\sum x_n$ es
        condicionalmente convergente, es decir, que $\sum|x_n|$ diverge. Esta contradicción muestra
        que $\sum p_n$ diverge. Como además $p_n\ge0$ para todo $n$, sus sumas parciales forman una
        sucesión no decreciente; si la serie diverge, es porque estas sumas parciales no están
        acotadas, es decir, $\sum p_n=+\infty$.

        El mismo argumento, intercambiando los papeles de $p_n$ y $q_n$, muestra que $\sum q_n$
        también diverge a $+\infty$: si $\sum q_n$ convergiera, entonces $\sum(x_n+q_n)=\sum p_n$
        convergería (pues $x_n+q_n=p_n$), y de nuevo $\sum|x_n|=\sum p_n+\sum q_n$ convergería —
        la misma contradicción.

        Finalmente, si $(x_n)$ tuviera solo un número finito de términos estrictamente positivos,
        entonces $p_n$ sería distinto de cero solo para finitos índices, y la serie $\sum p_n$ sería,
        en esencia, una suma finita, que converge trivialmente — contradiciendo lo que acabamos de
        probar. Así, $(x_n)$ tiene infinitos términos estrictamente positivos, es decir, la sucesión
        $(p_n')$ es infinita. El mismo argumento aplicado a $q_n$ muestra que $(x_n)$ tiene infinitos
        términos estrictamente negativos.
    \end{proof}

\end{document}